%! AP Statistics — Unit 2, Lesson 2.6 — Slides
\documentclass[aspectratio=169, 11pt]{beamer}
\usepackage{apstats-beamer}

\begin{document}

% ── TITLE SLIDE ───────────────────────────────────────────────────────────────
{
\setbeamercolor{background canvas}{bg=navy}
\begin{frame}[plain]
  \vspace{2.8cm}\hspace{0.55cm}%
  \begin{minipage}{0.9\textwidth}
    {\color{goldacc}\bfseries\small LESSON 2.6}\\[0.5cm]
    {\color{white}\bfseries\LARGE Linear Regression Models}\\[1.2cm]
    {\color{skymid}\small AP Statistics~$\cdot$~Unit 2: Exploring Two-Variable Data}
  \end{minipage}
\end{frame}
}

% ── LEARNING TARGETS ─────────────────────────────────────────────────────────
\begin{frame}[t, plain]
  \navyheader{Today's Learning Targets}
  \vspace{0.3cm}
  By the end of class, I will be able to\dots
  \begin{itemize}
    \item \textbf{write} the LSRL equation using the names of the variables.
    \item \textbf{interpret the slope} in context (direction, magnitude, units).
    \item \textbf{interpret the $y$-intercept} in context and assess whether it is meaningful.
    \item \textbf{make predictions} using the regression equation and distinguish
          interpolation from extrapolation.
  \end{itemize}
  \vspace{0.4cm}
  \textcolor{goldacc}{\small Standards: DAT-1.E [Skill 2.B] $\cdot$ DAT-1.F [Skill 2.B]}
\end{frame}

% ── HOOK ─────────────────────────────────────────────────────────────────────
\begin{frame}[t, plain]
  \navyheader{Hook --- Can You Predict a Shoe Size?}
  \vspace{0.3cm}
  \textit{The scatterplot shows height (inches) vs.\ shoe size for 30 students.}
  \vspace{0.3cm}
  \begin{columns}
    \column{0.52\textwidth}
      \textbf{Turn and talk:}
      \begin{itemize}
        \item If a student is 68 inches tall, what shoe size would you predict?
        \item How did you decide?
        \item Could you give a specific number using the scatterplot alone?
      \end{itemize}
    \column{0.44\textwidth}
      \textbf{Today's goal:}\\[6pt]
      Replace informal ``eyeballing'' with a \textbf{precise equation} that lets anyone
      compute the same prediction.
  \end{columns}
\end{frame}

% ── REGRESSION EQUATION ───────────────────────────────────────────────────────
\begin{frame}[t, plain]
  \navyheader{The Regression Equation}
  \vspace{0.3cm}
  \sectionlabel[goldacc]{DAT-1.E}
  \vspace{0.2cm}
  \[
    \hat{y} = a + bx
  \]
  \begin{columns}
    \column{0.48\textwidth}
      \begin{itemize}
        \item $b$ = \textbf{slope}: change in $\hat{y}$ per one-unit increase in $x$
        \item $a$ = \textbf{$y$-intercept}: predicted $\hat{y}$ when $x = 0$
        \item $\hat{y}$ = \textbf{predicted value} (not the actual $y$)
      \end{itemize}
    \column{0.48\textwidth}
      \textbf{AP Rule:} Use variable names, not $x$ and $y$.\\[8pt]
      \textcolor{navy}{$\widehat{\text{price}} = 28.4 - 1.85(\text{age})$}\\[4pt]
      \small\textit{Used car age (yr) vs.\ price (\$1000s)}
  \end{columns}
\end{frame}

% ── INTERPRETING THE SLOPE ────────────────────────────────────────────────────
\begin{frame}[t, plain]
  \navyheader{Interpreting the Slope}
  \vspace{0.25cm}
  \sectionlabel[goldacc]{DAT-1.E}
  \vspace{0.15cm}
  \textbf{Template:} ``For each additional [one unit of $x$], the predicted [y] [increases /
  decreases] by approximately [\,|b|\,].''
  \vspace{0.3cm}
  \begin{columns}
    \column{0.5\textwidth}
      \textbf{Car example:} $b = -1.85$\\[8pt]
      \textit{``For each additional year of age, the predicted resale price decreases by
      approximately \$1{,}850.''}
    \column{0.46\textwidth}
      \textbf{Three required pieces:}
      \begin{enumerate}
        \item One-unit increase in $x$ (with context)
        \item Direction (increases or decreases)
        \item Magnitude with units
      \end{enumerate}
      \vspace{4pt}
      \textcolor{redacc}{\small Missing any one $\Rightarrow$ partial credit at best.}
  \end{columns}
\end{frame}

% ── INTERPRETING THE Y-INTERCEPT ─────────────────────────────────────────────
\begin{frame}[t, plain]
  \navyheader{Interpreting the $y$-Intercept}
  \vspace{0.25cm}
  \sectionlabel[goldacc]{DAT-1.E}
  \vspace{0.15cm}
  \textbf{Template:} ``When [x-variable] $= 0$, the predicted [y-variable] is approximately [$a$].''
  \vspace{0.3cm}
  \begin{columns}
    \column{0.5\textwidth}
      \textbf{Car example:} $a = 28.4$\\[8pt]
      \textit{``When a car's age is 0 years (brand new), the predicted resale price is
      approximately \$28{,}400.''}\\[8pt]
      \textbf{Is it meaningful?} Yes --- age 0 is plausible.
    \column{0.46\textwidth}
      \textbf{When is it \emph{not} meaningful?}
      \begin{itemize}
        \item $x = 0$ is outside the data range, \textbf{OR}
        \item $x = 0$ is physically impossible
      \end{itemize}
      \vspace{6pt}
      \textcolor{redacc}{\small Always check before you interpret!}
  \end{columns}
\end{frame}

% ── PREDICTIONS ───────────────────────────────────────────────────────────────
\begin{frame}[t, plain]
  \navyheader{Making Predictions: Interpolation vs.\ Extrapolation}
  \vspace{0.25cm}
  \sectionlabel[goldacc]{DAT-1.F}
  \vspace{0.15cm}
  \begin{columns}
    \column{0.5\textwidth}
      \textbf{Predict} price of a 5-year-old car:\\[6pt]
      $\widehat{\text{price}} = 28.4 - 1.85(5) = 19.15$ (\$1000s)\\[8pt]
      Data ranged from \textbf{age 1 to 10 years}.\\[4pt]
      $5 \in [1, 10]$ $\Rightarrow$ \textbf{interpolation} $\Rightarrow$ reliable.
    \column{0.46\textwidth}
      \textbf{Predict} price of a 20-year-old car:\\[6pt]
      $\widehat{\text{price}} = 28.4 - 1.85(20) = -8.6$ (\$1000s)\\[8pt]
      $20 > 10$ $\Rightarrow$ \textbf{extrapolation} $\Rightarrow$ unreliable.\\[4pt]
      \textcolor{redacc}{\small A negative price is nonsensical ---
      the model breaks down outside the data range.}
  \end{columns}
\end{frame}

% ── LSRL THROUGH MEAN POINT ───────────────────────────────────────────────────
\begin{frame}[t, plain]
  \navyheader{Key Property: The LSRL Passes Through $(\bar{x},\, \bar{y})$}
  \vspace{0.3cm}
  \begin{columns}
    \column{0.55\textwidth}
      For the car data: $\bar{x} = 5.5$ yr, $\bar{y} = 18.225$ (\$1000s).\\[10pt]
      Verify: $28.4 - 1.85(5.5) = 28.4 - 10.175 = \mathbf{18.225}$ \\[6pt]
      \checkmark\ The line passes through $(\bar{x},\, \bar{y})$.
    \column{0.4\textwidth}
      \textbf{Why it matters:}\\[6pt]
      The mean point is always on the line. This is useful for checking your work and
      understanding why the line ``balances'' the data.
  \end{columns}
\end{frame}

% ── WRAP-UP ───────────────────────────────────────────────────────────────────
\begin{frame}[t, plain]
  \navyheader{Wrap-Up: What to Remember}
  \vspace{0.3cm}
  \begin{columns}
    \column{0.48\textwidth}
      \textbf{The regression equation:}
      \begin{itemize}
        \item Write with variable names (not $x$, $y$).
        \item Slope = predicted change in $\hat{y}$ per 1-unit $\uparrow$ in $x$.
        \item Intercept = predicted $\hat{y}$ when $x = 0$.
        \item Intercept: check if meaningful before interpreting.
      \end{itemize}
    \column{0.48\textwidth}
      \textbf{Predictions:}
      \begin{itemize}
        \item Interpolation: $x$ inside data range $\Rightarrow$ reliable.
        \item Extrapolation: $x$ outside data range $\Rightarrow$ unreliable, flag it.
        \item LSRL always passes through $(\bar{x},\, \bar{y})$.
      \end{itemize}
      \vspace{6pt}
      \textcolor{goldacc}{\small Next: Lesson 2.7 --- Residuals}
  \end{columns}
\end{frame}

\end{document}
